\section{Normalizador \(899\), unidad de Pell y estructura aritm\'etica del modo \(30\)}

El normalizador posee la factorizaci\'on aritm\'etica

\begin{equation}
899=30^2-1=29\cdot31.
\label{mab:rm:eq:899}
\end{equation}

El mismo modo \(30\) que genera los sectores \(90\) y \(120\) determina
simult\'aneamente el par de primos gemelos \(29,31\). Adem\'as,

\begin{equation}
\varepsilon_{30}=30+\sqrt{899},\qquad
\varepsilon_{30}^{-1}=30-\sqrt{899},\qquad
\log\varepsilon_{30}=\arcosh30.
\label{mab:rm:eq:pell30}
\end{equation}

Esta unidad de Pell determina la escala hiperb\'olica asociada al
normalizador y relaciona la fase dodecaf\'asica, el modo \(30\), los primos
gemelos y la geometr\'ia hiperb\'olica. Esta relaci\'on no identifica el
operador constitutivo del vac\'io con el renormalizador de duplicaci\'on; tal
identificaci\'on requerir\'ia un entrelazador expl\'icito.

La unidad act\'ua de forma integral mediante
\begin{equation}
M_{30}=\begin{pmatrix}30&899\\1&30\end{pmatrix},
\qquad \det M_{30}=1.
\label{mab:rm:eq:pell-matrix30}
\end{equation}
Sobre \(\Z[\sqrt{899}]\), esta matriz representa la multiplicaci\'on por
\(30+\sqrt{899}\). Sus autovalores son \(\varepsilon_{30}^{\pm1}\); por
tanto \(\log\varepsilon_{30}\) es la longitud de traslaci\'on de la transformaci\'on
hiperb\'olica asociada. El mismo entero central \(30\) reaparece en el vac\'io
porque \(90=3\cdot30\) y \(120=4\cdot30\), pero los dos operadores siguen
teniendo dominios distintos. La confluencia demostrada es la del modo y sus
completaciones; la identidad de din\'amicas no se postula.

